Data assimilation (DA) addresses the fundamental challenge of estimating unknown states or parameters in dynamical systems by optimally combining prior model predictions with noisy observations. This is typically accomplished through either variational methods, which minimize regularized objective functions, or Bayesian methods, which estimate posterior distributions -- usually via finite sample approximations \cite{evensen_data_2022}. The Bayesian approach is often preferred due to its natural framework for uncertainty quantification through posterior distribution estimation.

Within Bayesian DA, approaches such as ensemble Kalman filters and particle filters are arguably the most popular. Ensemble-based Kalman filters are computationally efficient but require Gaussianity assumptions and can only handle weakly nonlinear forward and observation models. Particle filters, conversely, impose virtually no restrictions but are computationally expensive due to the large number of samples required. This fundamental trade-off between computational tractability and modeling flexibility has motivated extensive research into alternative approaches.

One promising direction replaces expensive forward models with efficient neural network surrogates \cite{mucke_deep_2024, cheng_generalised_2022, zhong_reduced-order_2023, chen_reduced-order_2023}. While achieving substantial computational speed-ups, these deterministic approximations require specifying a priori model error distributions when incorporated into DA frameworks. This poses a fundamental challenge for uncertainty quantification in high-dimensional stochastic systems, where the true model error characteristics are unknown.

Recently, generative models have emerged as a compelling alternative, learning to directly sample from distributions of interest. In DA contexts, these models are trained on prior distributions and subsequently modified to incorporate observational constraints for posterior sampling. Two families of flow-based generative models dominate this landscape: stochastic interpolants (SIs) \cite{albergo_stochastic_2023, albergo_building_2023, chen_probabilistic_2024}, which transport samples between arbitrary distributions through a stochastic differential equation (SDE), and flow matching (FM) \cite{lipman_flow_2023}, which does so through a deterministic ordinary differential equation (ODE). Both have achieved remarkable success in computer vision and are increasingly applied to physical forecasting, yet methods for turning them into posterior samplers have so far been developed model by model, often tied to whether the generator is stochastic or deterministic.

\paragraph{A unified view.}
We argue that the distinction between SDE-based and ODE-based generators is, for the purpose of data assimilation, a choice of sampler rather than a difference in method. Both SIs and FM are special cases of an \emph{affine Gaussian probability path}, the common object underlying flows, diffusions, and interpolants \cite{albergo_stochastic_2023, lipman_flow_2023, ma_sit_2024, holderrieth_generator_2024}. This unification is itself established \cite{albergo_stochastic_2023, holderrieth_generator_2024}; our contribution is not the unification but its use as a single interface for posterior sampling, obtained by conditioning the path on the observation. Conditioning such a path on an observation tilts its target to the posterior while leaving the path structure intact, so the conditional score and the conditional velocity differ from their prior counterparts by the \emph{same} likelihood-score term -- scaled by one for the score and by the velocity--score duality coefficient for the velocity \cite{albergo_stochastic_2023, ma_sit_2024}. Because the conditioned path again admits a one-parameter family of equal-marginal dynamics with a freely chosen diffusion coefficient \cite{song_score-based_2021}, a single guidance term yields a whole family of posterior samplers. Three members are of practical interest: the stochastic-interpolant SDE driven by its native diffusion, the flow matching SDE obtained by lifting the probability-flow ODE, and a deterministic guided probability-flow ODE that steers flow matching directly without injecting noise \cite{yan_fig_2024, pokle_training-free_2024}. All three share one observation-interpolant likelihood and one multiplicative correction; they differ only in a diffusion coefficient. The paper presents the method as such, with the three samplers as worked instances of one framework.

\paragraph{Contributions:}
We introduce a posterior sampling method for data assimilation that applies to any flow-based generative model defined by an affine Gaussian probability path, encompassing both stochastic interpolants and flow matching without retraining. Our key contributions are:
\begin{itemize}
    \item We formulate data assimilation with flow-based generative models in a model-agnostic way by conditioning the affine Gaussian path on the observation. We show the conditioned path is again an affine Gaussian path whose score and velocity correct the prior ones by a single likelihood-score term, and we prove that the marginal-preserving family of this path generates exact posterior samples under ideal conditions.
    \item From this single family we obtain three posterior samplers -- the stochastic-interpolant SDE, the flow matching SDE, and a deterministic guided probability-flow ODE -- as the members with native, lifted, and zero diffusion. They share one guidance weight $\gweight_\tau=\vscoef_\tau+\tfrac12\gdiff_\tau^2$ that splits into a velocity--score duality term and a diffusion term, and we give pseudo-code for each.
    \item We introduce a practical likelihood-score estimator based on observation interpolation, derive its conditional moments for both model classes, and prove a closed-form multiplicative correction that maps the interpolant-likelihood score onto the standard Gaussian-surrogate (DPS-style) likelihood score \cite{chung_diffusion_2023, song_pseudoinverse_2023}; the correction is shared by all three samplers.
    \item We summarise the three samplers and the model-specific bookkeeping in two comparison tables (Tables~\ref{tab:samplers}--\ref{tab:si_vs_fm}), relate the construction to Doob's $h$-transform, and demonstrate it on a series of test cases.
\end{itemize}
